\documentclass[11pt]{article}
\usepackage[a4paper,margin=25mm]{geometry}
\usepackage{amsmath,amssymb,hyperref}
\title{Maximal heights on integer spheres}
\author{ziangni-sys}
\date{Conjecture 00000008342}
\begin{document}
\raggedright
\maketitle
\noindent\textbf{Result.} The asserted maximal-coordinate asymptotic $\sqrt{m/3}$ fails along every positive square $m=n^2$.

\paragraph{Scope and actual height.}
Both source versions describe $\mathrm{sph\_h}$ parenthetically as the maximal coordinate of integer sphere points. We use the corresponding extremal height
\[
S_m=\{(x,y,z)\in\mathbb Z^3:x^2+y^2+z^2=m\},\qquad
h(x,y,z)=\max(|x|,|y|,|z|),
\]
\[
H(m)=\sup\{h(p):p\in S_m\}.
\]
Only nonempty spheres are needed below. This is a maximum over all integer points, not a statistical mean or a typical value in a distribution. The source states the maximal-coordinate law without an exceptional restriction on square values of $m$.

\paragraph{Exact computation on an infinite family.}
For every integer $n\geq0$, the point $(n,0,0)$ belongs to $S_{n^2}$ and has height $n$. Conversely, every $(x,y,z)\in S_{n^2}$ satisfies
\[
x^2\leq n^2,\quad y^2\leq n^2,\quad z^2\leq n^2,
\]
because the other squared coordinates are nonnegative. Hence $|x|,|y|,|z|\leq n$ and $h(x,y,z)\leq n$. The height set is therefore nonempty and bounded, and its supremum is attained:
\[
\boxed{H(n^2)=n.}
\]
If ``maximal coordinate'' uses signed coordinates instead of absolute values, the maximum over all sphere points is still $n$: every signed coordinate is at most its absolute value, and $(n,0,0)$ attains $n$.

\paragraph{Asymptotic contradiction.}
For every $n>0$,
\[
\frac{H(n^2)}{\sqrt{n^2/3}}=\frac{n}{n/\sqrt3}=\sqrt3\ne1.
\]
The square parameters tend to infinity. Thus the ratio cannot tend to one even along this sequence of nonempty spheres, contradicting $H(m)\sim\sqrt{m/3}$. If the Chinese phrase is read as an additional factor $1/3$, the ratio to $(1/3)\sqrt{n^2/3}$ is instead $3\sqrt3$, also different from one. Both readings fail. No assertion about recurrence, exceptional-set density, or equidistribution estimates is required: the maximal-coordinate conjunct already fails.

\paragraph{Formal verification.}
The Lean project defines all integer triples satisfying the sphere equation, the coordinate height, and the actual real supremum of its image. It proves membership and height of the axis point, the bound for every sphere point, nonemptiness and boundedness of the height set, and the exact supremum. It proves that square parameters tend to infinity, computes the constant ratio and its actual limit, and rules out both proposed ratio-one limits. The results use the actual sphere and supremum, with no assumed height formula or numerical approximation. Lean~4.19.0 and a pinned Mathlib revision are included with axiom audits and reproduction instructions.
\end{document}
